\documentclass[10pt,a4paper]{amsart}
\usepackage[margin=19mm]{geometry}
\usepackage{amsmath,amssymb,mathtools}
\usepackage[scr=rsfs,cal=euler]{mathalfa}
\usepackage{xfrac}
\usepackage[unicode,hidelinks,bookmarks=false]{hyperref}
\newcommand{\ie}{i.e.\ }
\newcommand{\cf}{cf.\ }
\newcommand{\resp}{respectively\ }
\newcommand{\Subsection}{\subsection}
\providecommand{\nobreakdash}{\nobreak\hbox{-}}
\setlength{\emergencystretch}{2em}
\multlinegap0pt
\usepackage{bm}
\usepackage{bbm}
\usepackage{dsfont}
\usepackage{xspace}
\usepackage{cancel}
\usepackage{mathtools}
\usepackage{amsthm}
\makeatletter
\@ifundefined{thm}{\newtheorem{thm}{Theorem}}{}
\@ifundefined{example}{\newtheorem{example}[thm]{Example}}{}
\@ifundefined{lem}{\newtheorem{lem}[thm]{Lemma}}{}
\@ifundefined{notation}{\newtheorem{notation}[thm]{Notation}}{}
\@ifundefined{prop}{\newtheorem{prop}[thm]{Proposition}}{}
\makeatother
\newcommand{\glb}{\mbox{\rm glb}}
\newcommand{\qu}{\mathbb{Q}}
\newcommand{\re}{\mathbb{R}}
\newcommand{\vv}{\mathcal{Z}}
\title[Constructing Discontinuous but Locally Bounded Rational Func]{Constructing Discontinuous but Locally Bounded Rational Functions using \L ojasiewicz Inequalities}
\author{Adam Coffman, Yifei Pan}
\date{Source: arXiv:2606.01447v1 (CC BY 4.0)}
\begin{document}
\maketitle

\noindent\emph{Source and licence.} Constructing Discontinuous but Locally Bounded Rational Functions using \L ojasiewicz Inequalities. Version arXiv:2606.01447v1 (CC BY 4.0), distributed under the Creative Commons Attribution 4.0 International licence, as recorded in the arXiv licence metadata for this exact version and shown on the abstract page https://arxiv.org/abs/2606.01447v1. Full rights notice: licenses/arxiv-2606.01447-CC-BY-4.0.txt.\par\smallskip

\noindent\textit{Editorial scope.} Counted unit: the Thm whose source label is \texttt{thm3.3} in the article (source0.tex lines 562--575, source label \texttt{thm3.3}), delivered together with the article's own complete original proof of it (source0.tex lines 576--617, one \texttt{proof} environment of 42 lines). No mathematics is written, repaired or re-derived: statement and proof are the article's own text line for line. Required context retained uncounted: lem3.4 (lines 156--159); prop6.2 (lines 263--273); lem3.3 (lines 285--300); ex3.2 (lines 515--525); eq2 (lines 521--523); not4.5 (lines 528--530); thm3.2 (lines 537--546); eq4 (lines 549--552); eq7 (lines 549--552). Every definition, display and label of those lines is retained unchanged. Omitted from this package and not redistributed: 1--155, 160--262, 274--284, 301--514, 524--527, 531--536, 547--548, 553--561, 618--779. Nothing inside a retained excerpt is omitted or altered: every retained body is delivered exactly as the article's lines are written, so no omission receipt of any kind accompanies this package. Citations: the retained excerpts carry no citation command at all, so no bibliography entry of the article is transcribed and no reference is redistributed. Reference adaptation: none was needed and none was made; every reference inside the delivered unit resolves inside it, so no unresolved reference or citation is produced. Printed slips kept literal: no printed word, symbol, hypothesis or number of the counted statement or of its proof was altered. Layout and class: the article's own class is \texttt{amsart} with packages including amsmath, amssymb, amsthm, graphicx; the wrapper is the lane's pinned \texttt{amsart} editorial wrapper at 10pt on A4, which restates the theorem-like environments the retained text needs and adds the article's own macro definitions that the retained text needs (\texttt{\textbackslash glb}, \texttt{\textbackslash qu}, \texttt{\textbackslash re}, \texttt{\textbackslash vv}), with extra wrapper packages (bm, bbm, dsfont, xspace, cancel, mathtools). The delivered numbering is the unit's own. Assets: no retained span contains a figure environment, a graphics-inclusion command, a PDF-page inclusion, a table or a TikZ picture; no image file is copied into this package, the package declares no asset dependency and no image file is redistributed with it. Licence: version arXiv:2606.01447v1 of this article is distributed under the Creative Commons Attribution 4.0 International licence, as recorded in the arXiv licence metadata for this exact version and shown on the abstract page https://arxiv.org/abs/2606.01447v1. A full rights notice is delivered at licenses/arxiv-2606.01447-CC-BY-4.0.txt. Characters: every retained excerpt is pure ASCII, so the delivered engine, XeLaTeX with its default Latin Modern text font, reports no missing character.
\begin{lem}\label{lem3.4}
  Let $\Delta$ be any open set in $\re^n$.  If the rational function $P(\vec x)/Q(\vec x)$ is bounded on $\Delta$, then for any $\tau>0$, the function   $$S(\vec x)=\left\{\begin{array}{ll}{\displaystyle{\frac{P(\vec x)}{|Q(\vec x)|^{1-\tau}}}} & \mbox{\rm{for} $\vec x\in\Delta\setminus\vv(Q)$}\\ & \\0& \mbox{\rm{for} $\vec x\in\Delta\cap\vv(Q)$}\end{array}\right.$$
 is continuous on $\Delta$.
\end{lem}

\begin{prop}[\L ojasiewicz Gradient Inequality]\label{prop6.2}
  Given a polynomial $P$ and a point $\vec x_0\in\vv(P)\subseteq\re^n$, if $P\not\equiv0$ and $\vec\nabla P(\vec x_0)=\vec0$ then there is a radius $\rho>0$ so that $\vv(|\vec\nabla P|^2)\cap\overline B_\rho(\vec x_0)\subseteq\vv(P)$, and:
  \begin{itemize}
    \item {\rm{(\L5)}} This set ${\mathbf\Theta}\subseteq\re$ is non-empty:
    \begin{equation}\label{eq6.4}
  {\mathbf\Theta}=\left\{\theta>0:\exists r>0\ \ \exists C>0 : \  |\vec x-\vec x_0|<r\implies|P(\vec x)|^\theta\le C\left|\vec\nabla P(\vec x)\right|\right\}.\nonumber
 \end{equation}
   \item {\rm{(\L6)}} $\theta_0=\glb{\mathbf\Theta}$ satisfies $\theta_0\in\qu\cap\left[\frac12,1\right)$. %$\theta_0\in[\frac12,1)\cap\qu$.
    \item {\rm{(\L7)}} $\theta_0\in{\mathbf\Theta}$.
  \end{itemize}
\end{prop}

\begin{lem}\label{lem3.3}
  For an open ball $B_r(\vec x_0)=B_r$ and a rational function $P(\vec x)/Q(\vec x)$, if the function $s:B_r\to\re$ defined by
  $$s(\vec x)=\left\{\begin{array}{ll}{\displaystyle{\frac{P(\vec x)}{Q(\vec x)}}} & \mbox{\rm{for} $\vec x\in B_r\setminus\vv(Q)$}\\ & \\0& \mbox{\rm{for} $\vec x\in B_r\cap\vv(Q)$} \end{array}\right.$$ is continuous on $B_r$,
%  
%  $Q(\vec0)=0$ is an isolated zero, and 
%  \begin{equation}\label{eq3.3}
%  \lim_{\vec x\to\vec0}\frac{P(\vec x)}{Q(\vec x)}=0
%  \end{equation}
  then there exists a rational number $\varphi>0$ so that
    $$s_\varphi(\vec x)=\left\{\begin{array}{ll}{\displaystyle{\frac{|P(\vec x)|}{|Q(\vec x)|^{1+\varphi}}}} & \mbox{\rm{for} $\vec x\in B_r\setminus\vv(Q)$}\\ & \\0& \mbox{\rm{for} $\vec x\in B_r\cap\vv(Q)$} \end{array}\right.$$ is continuous on $B_r$.
%
%$P(\vec x)/(Q(\vec x))^{1+\varphi}$ is continuous in a neighborhood of $\vec0$ when extended by 
%  \begin{equation}\label{eq3.4}
%  \lim_{\vec x\to\vec0}\frac{P(\vec x)}{(Q(\vec x))^{1+\varphi}}=0.
%  \end{equation}
\end{lem}

\begin{example}\label{ex3.2}
  Start with a polynomial $P$ with critical point $\vec x_0$ as in Proposition \ref{prop6.2}, so $P(\vec x_0)=|\vec\nabla P(\vec x_0)|=0$.  If $P$ is non-constant then the squared norm $|\vec\nabla P(\vec x)|^2$ is a polynomial $\not\equiv0$.  Take any particular number $\theta$ in the set $\mathbf\Theta$ from statement (\L5) from Proposition \ref{prop6.2}, so there is some ball $B_r(\vec x_0)$, where $\vv(|\vec\nabla P|^2)\cap B_r\subseteq\vv(P)\cap B_r$.  For $\vec x\in B_r\setminus\vv(|\vec\nabla P|^2)$ and the bound $C$ corresponding to $\theta$ from Proposition \ref{prop6.2},% the following function is bounded by $C^2$,
\begin{equation*}%\label{eq6.1}
  0\le\frac{|P(\vec x)|^{2\theta}}{\left|\vec\nabla P(\vec x)\right|^2}\le C^2.
\end{equation*}
If we also assume there exist positive integers $p$, $q$, so that $p/q=\theta$, then raising to the power of $q$ gives:
\begin{equation}\label{eq2}
  0\le\frac{(P(\vec x))^{2p}}{\left|\vec\nabla P(\vec x)\right|^{2q}}\le C^{2q}.
\end{equation}
The rational function (\ref{eq2}) so constructed is bounded on $B_r$; the intersection of its indeterminacy set and $B_r$ is $\vv(|\vec\nabla P|^2)\cap B_r$.  
\end{example}

\begin{notation}\label{not4.5}
  Given $P$, $\vec x_0$, $p$, $q$ as in the construction of Example \ref{ex3.2}, the numerator and denominator of the expression (\ref{eq2}) each have a unique factorization (up to reordering and constant multiples, and not depending on $r$) into irreducible polynomials.  Let $P^{2p}=F\cdot N$ and $\left|\vec\nabla P(\vec x)\right|^{2q}=F\cdot D$, so that $F$ is a product of all of the common factors.  If there are no common factors, then let $F=1$; if   $P^{2p}$ has $\left|\vec\nabla P(\vec x)\right|^{2q}$ as a factor, then let $D=1$.  Canceling $F$ from the numerator and denominator of the rational function (\ref{eq2}) gives  another rational function $\displaystyle{\frac ND}$, the {\underline{reduced}} form of $\displaystyle{\frac{(P(\vec x))^{2p}}{\left|\vec\nabla P(\vec x)\right|^{2q}}}$. 
\end{notation}

\begin{thm}\label{thm3.2}
  Given a polynomial $P\not\equiv0$ with $P(\vec x_0)=|\vec\nabla P(\vec x_0)|=0$, positive integers $p$, $q$ so that  $p/q=\theta\in{\mathbf\Theta}$, and $r$, $C$, $N$ and $D$ as constructed in Example {\rm{\ref{ex3.2}}} and Notation {\rm{\ref{not4.5}}}, the rational function $N/D$ is bounded on $B_r$ and $\vv(D)\cap B_r$ has dimension $\le n-2$. 
  %For the following statements {\rm(1)} and {\rm(2)}, {\rm(1)}$\implies${\rm(2)}, and conversely, if $\vec x_0$ is an isolated zero of $D$, $\vv(D)\cap B_r=\{\vec x_0\}$, then {\rm(2)}$\implies${\rm(1)}.
%  \begin{enumerate}
%    \item $\theta>\theta_0$.
%      \item  The function 
%  $$f(\vec x)=\left\{\begin{array}{ll}N(\vec x)/D(\vec x) & \mbox{\rm{for} $\vec x\not\in\vv(D)$}\\ & \\0& \mbox{\rm{for} $\vec x\in\vv(D)$} \end{array}\right.$$
%  is continuous on some ball $B_{r^\prime}=\{\vec x:|\vec x|<r^\prime\}$.
%  \end{enumerate}
\end{thm}

\begin{eqnarray}
  \frac{N(\vec x)}{D(\vec x)}&=&\lim_{i\to\infty}\frac{N(\vec x_i)}{D(\vec x_i)}\label{eq4}\\
  &=&\lim_{i\to\infty}\frac{(P(\vec x_i))^{2p}}{\left|\vec\nabla P(\vec x_i)\right|^{2q}}\le C^{2q}.\label{eq7}
\end{eqnarray}

\begin{thm}\label{thm3.3}
  Given a polynomial $P\not\equiv0$ with $P(\vec x_0)=|\vec\nabla P(\vec x_0)|=0$, positive integers $p$, $q$ so that  $p/q=\theta\in{\mathbf\Theta}$, and $r$, $C$, $N$ and $D$ as constructed in Example {\rm{\ref{ex3.2}}} and Notation {\rm{\ref{not4.5}}}, define these functions $\re^n\to\re$,
    $$f(\vec x)=\left\{\begin{array}{ll}{\displaystyle{\frac{N(\vec x)}{D(\vec x)}}} & \mbox{\rm{for} $\vec x\not\in\vv(D)$}\\ & \\0& \mbox{\rm{for} $\vec x\in\vv(D)$} \end{array}\right.$$
and 
$$g(\vec x)=\left\{\begin{array}{ll}{\displaystyle{\frac{(P(\vec x))^{2p}}{|\vec\nabla P(\vec x)|^{2q}}}} & \mbox{\rm{for} $\vec\nabla P(\vec x)\ne\vec0$}\\ & \\0& \mbox{\rm{for} $\vec\nabla P(\vec x)=\vec0$} \end{array}\right..$$
The functions $f$ and $g$ are bounded on $B_{r}(\vec x_0)$, and further,
\begin{enumerate}
  \item If $p/q>\theta_0$, then $f(\vec x)$ and $g(\vec x)$ are equal and continuous on $B_{r^\prime}(\vec x_0)$ for some $r^\prime\le r$.

  \item If $p/q=\theta_0$, then $g(\vec x)$ is not continuous on any ball $B_{r^\prime}(\vec x_0)$.

  \item If $p/q=\theta_0$, and $\vec x_0$ is an isolated critical point of $P$, and $D(\vec x_0)=0$, then $f(x)$  is not continuous on any ball $B_{r^\prime}(\vec x_0)$.
  \end{enumerate}
\end{thm}

\begin{proof}
  The function $f$ is bounded on $B_r$ by Theorem \ref{thm3.2}.  By construction, $0\le g(\vec x)\le f(\vec x)$ on $\re^n$, with $g(\vec x)=f(\vec x)$ for $\vec x\in(\re^n\setminus\vv(|\vec\nabla P|^2))\cup\vv(D)$.   In the special case $\vv(D)=\vv(|\vec\nabla P|^2)$, $f$ and $g$ are equal everywhere.
  
  Recall $\theta_0$ is the smallest exponent from statement (\L7) of Proposition \ref{prop6.2}, and, from (\L6), let $p_0$, $q_0$ be positive integers so that $\displaystyle{\theta_0=\frac{p_0}{q_0}}$.  With the radius $r_0$ and $C_0$ corresponding to $\theta_0$ in the definition of $\mathbf\Theta$ from (\L5) in Proposition \ref{prop6.2}, for $\vec x\in B_{r_0}$, 
\begin{equation}\label{eq4.5}
  |P(\vec x)|^{\theta_0}\le C_0|\vec\nabla P(\vec x)|.
\end{equation}
The rational function $P^{2p_0}/|\vec\nabla P|^{2q_0}$ is bounded as in line (\ref{eq2}), and Lemma \ref{lem3.4} applies, so that for any $\tau>0$,
$$S(\vec x)=\left\{\begin{array}{ll}{\displaystyle{\frac{(P(\vec x))^{2p_0}}{(|\vec\nabla P(\vec x)|^{2q_0})^{1-\tau}}}} & \mbox{\rm{for} $\vec\nabla P(\vec x)\ne\vec0$}\\ & \\0& \mbox{\rm{for} $\vec\nabla P(\vec x)=\vec0$} \end{array}\right.$$
 is continuous on $B_{r_0}$.

 For statement (1), assuming $p/q=\theta>\theta_0$, let $\displaystyle{\tau=1-\frac{\theta_0}{\theta}>0}$, so that for $\vec\nabla P(\vec x)\ne\vec0$,
 \begin{equation}\label{eq8}
   \frac{(P(\vec x))^{2p_0}}{(|\vec\nabla P(\vec x)|^{2q_0})^{1-\tau}}=\frac{(P(\vec x))^{2p_0}}{|\vec\nabla P(\vec x)|^{2p_0q/p}}=\left(\frac{(P(\vec x))^{2p}}{|\vec\nabla P(\vec x)|^{2q}}\right)^{p_0/p}.
   \end{equation}
This shows $g(x)=(S(\vec x))^{p/p_0}$, so $g$ is continuous on $B_{r_0}$ and we set the claimed $r^\prime=\min\{r,r_0\}$.  We want to show that $f(\vec x)=(S(\vec x))^{p/p_0}$ on $B_{r_0}$.  If $D(\vec x)=0$ then $f(\vec x)=0=(S(\vec x))^{p/p_0}$.  For a point $\vec x\in B_{r_0}$ with $D(\vec x)\ne0$, some of the following limit steps are the same as (\ref{eq4}) and (\ref{eq7}), using a sequence $\vec x_i\in B_{r_0}\setminus\vv(|\vec\nabla P|^{2q})$, the continuity of $N/D$ and $S$ at $\vec x$, and equation (\ref{eq8}).
 \begin{equation}
  f(\vec x)=\frac{N(\vec x)}{D(\vec x)}=\lim_{i\to\infty}\frac{N(\vec x_i)}{D(\vec x_i)}=\lim_{i\to\infty}\frac{(P(\vec x_i))^{2p}}{\left|\vec\nabla P(\vec x_i)\right|^{2q}}=\lim_{i\to\infty}(S(\vec x_i))^{p/p_0}=(S(\vec x))^{p/p_0}.\nonumber
\end{equation}

For statement (2), we continue to use the above notation for $\displaystyle{\theta_0=\frac{p_0}{q_0}}$.  The following argument is the same as Case 2.\ from Theorem \ref{thm3.3}, using (\L7).  Suppose toward a contradiction that $g(\vec x)$ is continuous on some ball $B_{r^{\prime}}$.  Then Lemma \ref{lem3.3} applies: there is some rational $\varphi>0$ so that 
$$g_{\varphi}(\vec x)=\left\{\begin{array}{ll}{\displaystyle{\frac{(P(\vec x))^{2p_0}}{(|\vec\nabla P(\vec x)|^{2q_0})^{1+\varphi}}}} & \mbox{\rm{for} $\vec\nabla P(\vec x)\ne\vec0$}\\ & \\0& \mbox{\rm{for} $\vec\nabla P(\vec x)=\vec0$} \end{array}\right.$$
is continuous, and bounded by $C_\varphi$ on some possibly smaller ball $B_{r^{\prime\prime}}$.
It follows that for $\vec x\in B_{r^{\prime\prime}}\setminus\vv(|\vec\nabla P(\vec x)|^2)$, 
\begin{equation}\label{eq4.6}
|P(\vec x)|^{2p_0/(2q_0(1+\varphi))}\le C_\varphi^{1/(2q_0(1+\varphi))}|\vec\nabla P(\vec x)|.
\end{equation}
If $\vec\nabla P(\vec x)=\vec0$ then both sides of the inequality (\ref{eq4.6}) are $=0$ by the inequality (\ref{eq4.5}), so the inequality (\ref{eq4.6}) holds for all $\vec x\in B_{r^{\prime\prime}}$.  This contradicts the minimality of $\theta_0$.

For statement (3), if there is some $r^{\prime\prime}$ so that $\vv(|\vec\nabla P|^2)\cap B_{r^{\prime\prime}}=\{\vec x_0\}$, then $\vv(|\vec\nabla P|^2)\cap B_{r^{\prime\prime}}=\vv(D)\cap B_{r^{\prime\prime}}$ and so $f=g$ on $B_{r^{\prime\prime}}$.  Given any $B_{r^\prime}$, $g$ is not continuous on $B_{r^\prime}\cap B_{r^{\prime\prime}}$ by (2) and $f$ is not continuous on $B_{r^\prime}$. 
%
%For $\theta=p/q\ge\theta_0$ as in Example \ref{ex3.2}, and for $\vec x\in B_{r_0}\setminus\vv(|\vec\nabla P|^2)$,
%    \begin{eqnarray}
%        |P(\vec x)|^{2\theta_0q}&\le&C_0^{2q}\left|\vec\nabla P(\vec x)\right|^{2q}\nonumber\\
%        \implies |P(\vec x)|^{2p}&=&|P(\vec x)|^{2p-2\theta_0q}|P(\vec x)|^{2\theta_0q}\le|P(\vec x)|^{2p-2\theta_0q}C_0^{2q}\left|\vec\nabla P(\vec x)\right|^{2q}\nonumber\\
%        \implies 0\le\frac{(P(\vec x))^{2p}}{\left|\vec\nabla P(\vec x)\right|^{2q}}&\le&C_0^{2q}|P(\vec x)|^{2p-2\theta_0q}.\label{eq5}
%    \end{eqnarray}
%    The exponent $2p-2\theta_0q$ is non-negative by the assumption that $\frac pq\ge\theta_0$.  
%    
%    For $\theta=p/q$, $r$, and $C$ as in the above construction, the inequality (\ref{eq2}) holds for $\vec x\in B_r\setminus\vv(|\vec\nabla P|^2)$. Let $r_1=\min\{r,r_0\}$, Similarly, w
%       At points where $D\ne0$, the same limit calculation as (\ref{eq4}) shows that $N/D$ has the same bounds as (\ref{eq5}) on $B_{r_0}$.  If the exponent By the squeeze theorem and the construction of $f$, $$\lim_{\vec x\to\vec x_0}f(\vec x)=0.$$
\end{proof}

\end{document}
